\documentclass{situationist-preprint}

\title{Institutional Uptake as a Criterion}
\msid{AES-01-b}
\msfamily{AES-01}
\msdiscipline{Aesthetics}
\mstype{Companion manuscript}
\mskeywords{institutional theory; artworld; critical uptake; presentation; Wollheim's dilemma; definition of art}

\begin{document}
\maketitle

\begin{abstract}
AES-01 holds that art is what is presented as art within an artworld, subject to the artworld's critical uptake. We argue that the second condition is not an optional refinement but the part of the account that makes it a criterion. Presentation alone is too permissive, since anyone can offer anything; market uptake tracks collectibility, not status; and conferral by an authorized agent faces Wollheim's dilemma between arbitrariness and redundancy. Critical uptake avoids all three difficulties. It is public, it is defeasible, it includes adverse judgment, and it involves reasons without being reducible to them. We state the criterion, show that it classifies the standard cases correctly, and answer the objections that it is circular and that it makes art status depend on fashion.
\end{abstract}

\section{The problem}

Institutional accounts of art descend from \citet{Danto1964}, who argued that to see something as art requires an atmosphere of theory and a knowledge of the history of art, and from \citet{Dickie1974}, who defined a work of art as an artefact upon which some person or persons acting on behalf of the artworld have conferred the status of candidate for appreciation. Dickie's definition was widely criticized for leaving it obscure what the conferral consists in and who is entitled to perform it. Dickie later replaced conferral by a network of roles \citep{Dickie1984}, in which an artist creates something to be presented to an artworld public; but the revised account was then charged with circularity, since artist, artworld and public are each defined by reference to the others.

The account of AES-01 has two conditions: an object must be presented as art within an artworld, and it is subject to the artworld's critical uptake. We use the two notions as AES-15-a defines them. An object is presented as art if some participant in an artworld offers it to a public as a candidate for the kind of appreciation and criticism the artworld practices (Definition~2.1 of AES-15-a); a presented object receives critical uptake if the artworld's critical practice engages it as a candidate, so that it is interpreted, compared with other works and assessed by the standards the practice applies to artworks (Definition~2.2 of AES-15-a). Our question is why the second condition is needed, and what makes it a criterion rather than a restatement of the problem.

\section{Three inadequate criteria}

\begin{definition}\label{d:criterion}
A condition on art status is a \emph{criterion} if (i) whether an object meets it is in principle publicly ascertainable, (ii) it can fail to be met by an object that someone sincerely offers as art, and (iii) the fact that it is met is not settled merely by someone's decision that it is.
\end{definition}

Condition (i) excludes accounts that make art status depend on a private mental act. Condition (ii) excludes accounts on which offering suffices. Condition (iii) is the requirement whose violation drives Wollheim's objection.

\begin{proposition}\label{p:pres}
Presentation alone is not a criterion.
\end{proposition}

\begin{proof}
Presentation in the sense of Definition~2.1 of AES-15-a is met whenever some participant in an artworld offers an object as a candidate. Since any participant may offer any object, a sincere offer by a participant always meets it, and condition (ii) of Definition~\ref{d:criterion} fails.
\end{proof}

\begin{proposition}\label{p:market}
Market uptake is not a criterion of art status.
\end{proposition}

\begin{proof}
An object receives market uptake when it is bought and sold at prices that reflect its supposed standing. Market uptake fails condition (iii) in one direction, since a sale is settled by the decision of a buyer and a seller; and it is neither necessary nor sufficient for status. Many works that the artworld treats as works were never sold, and objects are sold at high prices as collectibles, relics or investments without being interpreted or assessed as works. AES-15-a draws the same distinction in its reply to the objection that the sale of a machine-generated image showed uptake.
\end{proof}

\begin{proposition}\label{p:confer}
Conferral by an authorized agent is not a criterion.
\end{proposition}

\begin{proof}
\citet{Wollheim1980} put the difficulty as a dilemma. Either the representatives of the artworld confer status for reasons, or they do not. If they do, the reasons themselves determine which objects are art, and the institutional apparatus is redundant. If they do not, the status they confer is arbitrary, and nothing explains why it should be valued. Conferral by decision meets condition (ii) but fails condition (iii) on the second horn and is idle on the first.
\end{proof}

\section{Critical uptake}

\begin{proposition}\label{p:uptake}
Critical uptake is a criterion in the sense of Definition~\ref{d:criterion}.
\end{proposition}

\begin{proof}
(i) Interpretation, comparison and assessment are public activities: they issue in reviews, catalogues, exhibitions, scholarship, teaching and the responses of other artists, and whether an object has been engaged in this way is a matter of record. (ii) An object can be offered by a participant and ignored by the practice, or engaged only as a curiosity, a provocation or an episode in the history of a technology, and then it has not received critical uptake. (iii) Uptake is not an act of decision by any single agent. It consists in a pattern of engagement by many participants, each of whom judges the object on grounds internal to the practice, and no one of whom can bring the pattern about by declaring it.
\end{proof}

The criterion escapes Wollheim's dilemma. Critical uptake involves reasons, since to interpret and assess an object is to give reasons about it. But the reasons are not reasons for the classification \enquote{art}; they are reasons for particular readings and evaluations. What the practice confers is not a verdict that the object is art, from which the reasons would be detachable, but a place among the objects about which its reasons are given. The first horn assumes that if there are reasons, there is a reason-based definition of art that makes the institution redundant. There is not: the reasons differ from work to work and from period to period, and what they have in common is only that they are given within one practice. The second horn assumes that without such a definition the status is arbitrary. It is not, since each act of engagement answers to standards that participants dispute and can get wrong.

Two features of the criterion deserve emphasis. First, uptake includes adverse judgment. A work that critics engage and find weak, derivative or meretricious has been taken up: it has been assessed as a work. Status and merit come apart, as they must if there is bad art. Second, uptake is defeasible and graded. An object may be engaged by a small part of the practice and ignored by the rest, and its status may then be contested; the account says that such cases are contested because the condition is partially met, and this is the right verdict about them.

\begin{objection}[Circularity]\label{o:circ}
Critical uptake is engagement by the artworld's critical practice, and the practice is identified as the one that engages artworks. The criterion presupposes what it defines.
\end{objection}

\begin{reply}
The practice is identified historically, not by the extension of \enquote{art}. It is a continuing practice of making, exhibiting, interpreting and assessing, with institutions, vocabularies and canons that can be traced through time. \citet{Levinson1979} made the parallel point that art status can be fixed by relation to the history of art-regard without a prior definition. A definition that appeals to an independently identifiable practice is not circular, even if its terms are learned together.
\end{reply}

\begin{objection}[Fashion]\label{o:fashion}
If uptake makes art, then art status depends on the vagaries of critical fashion, and a work could become art or cease to be art as reputations change.
\end{objection}

\begin{reply}
Reputation and status differ. Fashion moves reputations, which are matters of merit; uptake in our sense survives a fall in reputation, since adverse assessment is still assessment. An object that was once engaged as a work and is now neglected remains an object the practice has engaged. What the criterion allows is that an object offered and ignored may later be taken up, as happened with works rediscovered after their makers' deaths. That is a feature of the practice and not an artefact of the account.
\end{reply}

\section{Conclusion}

Presentation, market uptake and conferral each fail as criteria of art status. Critical uptake succeeds: it is public, it can fail to be met by a sincere offer, and it is not settled by anyone's decision. It answers Wollheim's dilemma by locating reasons inside the practice without making them a definition that renders the practice redundant. The second condition of AES-01 is therefore what turns the institutional account from a description of how art is offered into a criterion of what art is.

\begin{thebibliography}{99}\small
\bibitem[Danto(1964)]{Danto1964} Danto, A.~C. (1964). The artworld. \emph{The Journal of Philosophy}, 61(19), 571--584.
\bibitem[Dickie(1974)]{Dickie1974} Dickie, G. (1974). \emph{Art and the Aesthetic: An Institutional Analysis}. Cornell University Press.
\bibitem[Dickie(1984)]{Dickie1984} Dickie, G. (1984). \emph{The Art Circle: A Theory of Art}. Haven Publications.
\bibitem[Levinson(1979)]{Levinson1979} Levinson, J. (1979). Defining art historically. \emph{The British Journal of Aesthetics}, 19(3), 232--250.
\bibitem[Wollheim(1980)]{Wollheim1980} Wollheim, R. (1980). The institutional theory of art. In \emph{Art and Its Objects}, 2nd edition. Cambridge University Press.
\end{thebibliography}

\end{document}
